% ===== PRÉAMBULE CANONIQUE — à coller VERBATIM en tête de CHAQUE .tex =====
\documentclass[11pt,a4paper]{article}
\usepackage[utf8]{inputenc}\usepackage[T1]{fontenc}\usepackage{lmodern}
\usepackage[french]{babel}
\usepackage{amsmath,amssymb,mathtools}\usepackage{icomma}
\usepackage[version=4]{mhchem}          % \ce{...} : équations & formules
\usepackage{siunitx}\sisetup{locale=FR} % \qty{}{}, \si{}, \ang{}
\usepackage{chemfig}                    % \chemfig{...} : structures développées
\usepackage{xcolor}\usepackage{colortbl}
\usepackage{tikz}\usepackage{pgfplots}\pgfplotsset{compat=1.18}
\usepackage[most]{tcolorbox}\usepackage{enumitem}\usepackage{array,booktabs}
\usepackage[a4paper,margin=2.2cm]{geometry}
\usetikzlibrary{arrows.meta,calc,angles,quotes,positioning,patterns,backgrounds,shapes.geometric}
\definecolor{primary}{RGB}{222,49,99}\definecolor{primarydark}{RGB}{150,22,55}
\definecolor{defcol}{RGB}{0,114,178}\definecolor{methcol}{RGB}{52,92,148}
\definecolor{excol}{RGB}{0,135,90}\definecolor{attncol}{RGB}{200,40,55}
\tcbset{boxbase/.style={enhanced,breakable,boxrule=0.9pt,arc=2.5pt,
  left=3mm,right=3mm,top=2.3mm,bottom=2.3mm,fonttitle=\bfseries,coltitle=white}}
% --- boîtes TUTORIEL (noms EXACTS, ne pas renommer) ---
\newtcolorbox{cle}[1][]{boxbase,colback=defcol!6,colframe=defcol,
  title={\textsc{À retenir}\ifx\relax#1\relax\relax\else\textmd{ : \itshape #1}\fi}}
\newtcolorbox{methode}[1][]{boxbase,colback=methcol!7,colframe=methcol,sharp corners,
  title={\textsc{Méthode}\ifx\relax#1\relax\relax\else\textmd{ --- \itshape #1}\fi}}
\newtcolorbox{exemplebox}[1][]{boxbase,colback=excol!5,colframe=excol,
  title={\textsc{Exemple}\ifx\relax#1\relax\relax\else\textmd{ : \itshape #1}\fi}}
\newtcolorbox{piege}[1][]{boxbase,colback=attncol!6,colframe=attncol,
  title={\textsc{Piège}\ifx\relax#1\relax\relax\else\textmd{ : \itshape #1}\fi}}
\newtcolorbox{remarque}[1][]{boxbase,colback=black!4,colframe=black!45,
  title={\textsc{Remarque}\ifx\relax#1\relax\relax\else\textmd{ : \itshape #1}\fi}}
% --- boîtes EXERCICE / CORRIGÉ (noms EXACTS, auto-comptées) ---
\newtcolorbox[auto counter]{exo}[1][]{boxbase,colback=primary!4,colframe=primary,
  title={\textsc{Exercice}~\thetcbcounter\ifx\relax#1\relax\relax\else\textmd{ --- \itshape #1}\fi}}
\newtcolorbox[auto counter]{corrige}[1][]{boxbase,colback=excol!4,colframe=excol,
  title={\textsc{Corrigé}~\thetcbcounter\ifx\relax#1\relax\relax\else\textmd{ --- \itshape #1}\fi}}
\newcommand{\R}{\mathbb{R}}\newcommand{\N}{\mathbb{N}}\newcommand{\Z}{\mathbb{Z}}\newcommand{\ds}{\displaystyle}
% (un \usepackage{fancyhdr}+en-tête est possible ; il est ignoré côté web)
% =========================================================================
\begin{document}

\begin{corrige}[métal à valence variable et anion simple]
\begin{enumerate}[label=\alph*)]
  \item \ce{FeCl3} : croiser $3$ et $1$ ; pas de parenthèses (\ce{Cl} monoatomique).
  \item \ce{CuO} : croiser $2$ et $2$ donne \ce{Cu2O2}, à \emph{simplifier} en \ce{CuO}.
  \item \ce{SnO2} : croiser $4$ et $2$ donne \ce{Sn2O4}, à simplifier en \ce{SnO2}.
\end{enumerate}
\end{corrige}

\begin{corrige}[métal et anion polyatomique]
\begin{enumerate}[label=\alph*)]
  \item \ce{CaCO3} : valences $2$ et $2 \Rightarrow \ce{Ca2(CO3)2}$, simplifié en \ce{CaCO3}.
  \item \ce{Al2(SO4)3} : croiser $3$ et $2$ ; \ce{SO4} parenthésé (présent 3 fois).
  \item \ce{Sr(IO3)2} : \ce{Sr^2+} (valence 2) et \ce{IO3-} (valence 1) $\Rightarrow$ deux iodates, parenthésés.
\end{enumerate}
\end{corrige}

\begin{corrige}[sels d'ammonium]
\begin{enumerate}[label=\alph*)]
  \item \ce{(NH4)2SO3} : deux \ce{NH4+} pour un \ce{SO3^2-}.
  \item \ce{(NH4)3PO4} : trois \ce{NH4+} pour un \ce{PO4^3-}.
  \item \ce{(NH4)2CO3} : deux \ce{NH4+} pour un \ce{CO3^2-} (carbonate d'ammonium).
\end{enumerate}
\end{corrige}

\begin{corrige}[du nom complet à la formule]
On identifie cation, anion et leurs valences, puis on croise.
\begin{enumerate}[label=\alph*)]
  \item sulfate d'aluminium : \ce{Al^3+} + \ce{SO4^2-} $\Rightarrow$ \ce{Al2(SO4)3}.
  \item nitrate de calcium : \ce{Ca^2+} + \ce{NO3-} $\Rightarrow$ \ce{Ca(NO3)2}.
  \item phosphate de sodium : \ce{Na+} + \ce{PO4^3-} $\Rightarrow$ \ce{Na3PO4}.
\end{enumerate}
\end{corrige}

\begin{corrige}[distinguer des anions voisins]
La racine et le suffixe fixent l'anion : sulf\emph{ite} \ce{SO3^2-}, sulf\emph{ate} \ce{SO4^2-},
\emph{thio}sulfate \ce{S2O3^2-} (tous de valence 2).
\begin{enumerate}[label=\alph*)]
  \item sulfite de magnésium : \ce{Mg^2+} + \ce{SO3^2-} $\Rightarrow$ \ce{MgSO3}.
  \item sulfate de magnésium : \ce{Mg^2+} + \ce{SO4^2-} $\Rightarrow$ \ce{MgSO4}.
  \item thiosulfate de sodium : \ce{Na+} + \ce{S2O3^2-} $\Rightarrow$ \ce{Na2S2O3}.
\end{enumerate}
\end{corrige}

\end{document}
